\documentclass[10pt, aspectratio=169]{beamer}
\usepackage{../beamer_theme_408}
\title{408考研零基础 --- 四科串联总结}
\subtitle{5-1 从输入URL到页面显示}
\author{}
\date{}
\begin{document}

\begin{frame}{本节目标}
    \begin{enumerate}
        \item 串联计算机网络、操作系统、计算机组成原理、数据结构四科知识
        \item 理解从输入URL到页面显示的完整过程
        \item 掌握各科在完整流程中的角色与协作关系
    \end{enumerate}
\end{frame}

\begin{frame}{课程导入：四科串联的经典场景}
    \begin{columns}[T]
        \column{0.5\textwidth}
        \textbf{前四个模块的知识回顾}
        \begin{itemize}
            \item \textbf{计网}：数据如何从源端传输到目的端
            \item \textbf{OS}：如何管理进程、内存、文件、I/O
            \item \textbf{计组}：CPU如何执行指令、存储器层次
            \item \textbf{数据结构}：数据如何高效组织和查找
        \end{itemize}
        \vspace{0.5em}
        \alert{各科并非孤立的，在实际场景中协同工作}
        \column{0.5\textwidth}
        \begin{center}
        \begin{tikzpicture}[>=stealth]
            \node[draw=408blue,fill=408blue!5,rounded corners,minimum width=3cm,minimum height=0.8cm] (net) {计算机网络};
            \node[draw=408blue,fill=408blue!5,rounded corners,minimum width=3cm,minimum height=0.8cm,below=0.3of net] (os) {操作系统};
            \node[draw=408blue,fill=408blue!5,rounded corners,minimum width=3cm,minimum height=0.8cm,below=0.3of os] (arch) {计算机组成原理};
            \node[draw=408blue,fill=408blue!5,rounded corners,minimum width=3cm,minimum height=0.8cm,below=0.3of arch] (ds) {数据结构};
            \draw[->,thick,408orange] (net.east) -- ++(0.5,0) -- ++(0,-0.75) -- (os.east);
            \draw[->,thick,408orange] (os.east) -- ++(0.5,0) -- ++(0,-0.75) -- (arch.east);
            \draw[->,thick,408orange] (arch.east) -- ++(0.5,0) -- ++(0,-0.75) -- (ds.east);
            \node[right=0.5of net,text=408orange,font=\small] {传输};
            \node[right=0.5of os,text=408orange,font=\small] {调度};
            \node[right=0.5of arch,text=408orange,font=\small] {执行};
            \node[right=0.5of ds,text=408orange,font=\small] {组织};
        \end{tikzpicture}
        \end{center}
    \end{columns}
\end{frame}

\begin{frame}{完整流程概览}
    \begin{center}
    \begin{tikzpicture}[node distance=3mm, auto]
        \tikzset{
            box/.style={rectangle, draw=primary, fill=white!90!primary, rounded corners, 
                        text centered, minimum width=20mm, minimum height=8mm, font=\small},
            arrow/.style={->, thick, primary}
        }
        \node[box] (url) {输入URL};
        \node[box, right=of url] (dns) {DNS解析};
        \node[box, right=of dns] (http) {HTTP请求};
        \node[box, right=of http] (tcp) {TCP连接};
        \node[box, right=of tcp] (ip) {IP路由};
        \node[box, below=of ip] (arp) {ARP-MAC};
        \node[box, below=of arp] (frame) {封装成帧};
        \node[box, right=of ip] (serve) {服务器处理};
        \node[box, below=of serve] (resp) {HTTP响应};
        \node[box, right=of resp] (render) {浏览器渲染};

        \draw[arrow] (url) -- (dns);
        \draw[arrow] (dns) -- (http);
        \draw[arrow] (http) -- (tcp);
        \draw[arrow] (tcp) -- (ip);
        \draw[arrow] (ip) -- (arp);
        \draw[arrow] (arp) -- (frame);
        \draw[arrow] (frame) -| (serve);
        \draw[arrow] (serve) -- (resp);
        \draw[arrow] (resp) -| (render);
    \end{tikzpicture}
    \end{center}
    \vspace{0.5em}
    \centering\textcolor{secondary}{贯穿科目：计算机网络~$\cdot$~操作系统~$\cdot$~计组~$\cdot$~数据结构}
\end{frame}

\begin{frame}{URL解析：浏览器第一步}
    \begin{block}{URL分解}
        \texttt{http://www.example.com:80/index.html}
    \end{block}
    \vspace{0.3em}
    \begin{columns}[T]
        \column{0.5\textwidth}
        \textbf{URL组成部分}
        \begin{itemize}
            \item \textbf{协议}：http / https（默认80/443端口）
            \item \textbf{域名}：www.example.com（需DNS解析）
            \item \textbf{端口}：:80（可省略使用默认值）
            \item \textbf{路径}：/index.html（资源定位）
        \end{itemize}
        \column{0.5\textwidth}
        \textbf{浏览器预处理}
        \begin{itemize}
            \item HSTS预加载列表检查（强制HTTPS）
            \item 检查是否有缓存的重定向（301/302）
            \item 解析域名前查本地hosts文件
        \end{itemize}
    \end{columns}
    \vspace{0.3em}
    \alert{URL解析是后续一切操作的起点，必须将域名解析为IP地址才能建立连接}
\end{frame}

\begin{frame}{计网：DNS域名解析}
    \begin{block}{DNS（域名系统）}
        \begin{itemize}
            \item 将 \texttt{www.example.com} 解析为 IP 地址
            \item 查询顺序：浏览器缓存 $\to$ 操作系统缓存 $\to$ 本地DNS服务器 $\to$ 根/顶级/权威DNS
        \end{itemize}
    \end{block}
    \begin{block}{涉及的数据结构}
        \begin{itemize}
            \item \textbf{DNS缓存：哈希表} --- O(1) 查找已解析的域名
            \item \textbf{域名层级：树形结构} --- 根、顶级域、二级域、子域
        \end{itemize}
    \end{block}
    \vspace{0.3em}
    \alert{DNS使用UDP协议（53端口），响应中包含TTL控制缓存有效期}
\end{frame}

\begin{frame}{DNS解析：递归查询 vs 迭代查询}
    \begin{columns}[T]
        \column{0.48\textwidth}
        \textbf{递归查询}
        \begin{itemize}
            \item 客户端向本地DNS服务器发起查询
            \item 本地DNS服务器负责完成整个查询过程
            \item 最终返回最终结果给客户端
            \item \textcolor{blue}{查询负担集中在本地DNS服务器}
        \end{itemize}
        \vspace{0.3em}
        \textbf{流程图}：\\
        客户端 $\to$ 本地DNS（递归）$\to$ 根DNS\\
        \hspace{4.5em}$\to$ 顶级DNS $\to$ 权威DNS
        \column{0.48\textwidth}
        \textbf{迭代查询}
        \begin{itemize}
            \item 本地DNS服务器依次向各DNS服务器发起查询
            \item 每个DNS服务器返回下一步的服务器地址
            \item 本地DNS服务器"迭代"逐级查询
            \item \textcolor{blue}{减轻各级DNS的负担}
        \end{itemize}
        \vspace{0.3em}
        \textbf{实际场景}：\\
        客户端$\to$本地DNS（递归）\\
        本地DNS$\to$根DNS（迭代）\\
        本地DNS$\to$顶级DNS（迭代）\\
        本地DNS$\to$权威DNS（迭代）
    \end{columns}
    \vspace{0.3em}
    \alert{考试重点：递归查询是"谁帮我查完"，迭代查询是"告诉我下一步找谁"}
\end{frame}

\begin{frame}{计网：HTTP请求 --- TCP三次握手}
    \begin{block}{HTTP请求}
        \begin{itemize}
            \item 浏览器构建HTTP请求报文（请求行、首部、空行、主体）
            \item HTTP/1.1默认持久连接，HTTP/2支持多路复用
        \end{itemize}
    \end{block}
    \begin{block}{TCP三次握手}
        \begin{enumerate}
            \item 客户端 $\to$ 服务器：\texttt{SYN=1, seq=x}
            \item 服务器 $\to$ 客户端：\texttt{SYN=1, ACK=1, seq=y, ack=x+1}
            \item 客户端 $\to$ 服务器：\texttt{ACK=1, seq=x+1, ack=y+1}
        \end{enumerate}
    \end{block}
    \vspace{0.3em}
    \alert{TCP头部至少20字节，包含源/目的端口、序号、确认号、窗口大小等字段}
\end{frame}

\begin{frame}{TCP三次握手与HTTP的关系（时序）}
    \begin{center}
    \begin{tikzpicture}[>=stealth, node distance=1.2cm]
        \node (client) at (0,0) {客户端};
        \node (server) at (6,0) {服务器};
        \draw[dashed] (client) -- (0,-5);
        \draw[dashed] (server) -- (6,-5);
        \draw[->,blue,thick] (0.2,-0.3) -- node[above,sloped]{\small SYN, seq=x} (5.8,-0.8);
        \draw[->,red,thick] (5.8,-1.2) -- node[above,sloped]{\small SYN+ACK, seq=y, ack=x+1} (0.2,-1.7);
        \draw[->,blue,thick] (0.2,-2.1) -- node[above,sloped]{\small ACK, seq=x+1, ack=y+1} (5.8,-2.6);
        \node at (3,-3.0) {\small \textbf{三次握手完成，连接建立}};
        \draw[->,blue,thick] (0.2,-3.7) -- node[above,sloped]{\small HTTP GET请求 (数据)} (5.8,-4.2);
        \node[right,font=\small] at (6.2,-4.0) {传输HTTP数据};
    \end{tikzpicture}
    \end{center}
    \vspace{0.3em}
    \alert{关键理解：TCP三次握手发生在HTTP数据发送\textbf{之前}，握手完成后才传输真正的HTTP请求}
\end{frame}

\begin{frame}{IP数据报转发过程}
    \begin{center}
    \begin{tikzpicture}[>=stealth, node distance=1cm]
        \node[draw,circle,minimum size=1cm,fill=blue!10] (host) {主机};
        \node[draw,rectangle,minimum width=1.5cm,minimum height=0.8cm,right=of host,fill=green!10] (r1) {路由器1};
        \node[draw,rectangle,minimum width=1.5cm,minimum height=0.8cm,right=of r1,fill=green!10] (r2) {路由器2};
        \node[draw,rectangle,minimum width=1.5cm,minimum height=0.8cm,right=of r2,fill=green!10] (r3) {路由器3};
        \node[draw,circle,minimum size=1cm,right=of r3,fill=red!10] (server) {服务器};
        \draw[->,thick] (host) -- node[above,font=\small] {查路由表} (r1);
        \draw[->,thick] (r1) -- node[above,font=\small] {下一跳} (r2);
        \draw[->,thick] (r2) -- node[above,font=\small] {下一跳} (r3);
        \draw[->,thick] (r3) -- node[above,font=\small] {到达} (server);
        \node[below=0.3of r2,font=\small] {\textcolor{408orange}{逐跳转发（Hop-by-Hop）}};
    \end{tikzpicture}
    \end{center}
    \vspace{0.3em}
    \begin{block}{路由过程}
        \begin{enumerate}
            \item 主机判断目的IP是否在同一子网
            \item 不在同一子网则发送给默认网关（路由器）
            \item 每个路由器查路由表做最长前缀匹配
            \item 逐跳转发直到到达目标网络
        \end{enumerate}
    \end{block}
    \alert{路由表查找失败 $\to$ 触发ICMP"目的不可达"报文}
\end{frame}

\begin{frame}{ARP协议：IP地址到MAC地址的映射}
    \begin{block}{为什么需要ARP？}
        \begin{itemize}
            \item IP数据报在网络层使用IP地址标识
            \item 实际传输时，数据链路层需要\textbf{MAC地址}来封装帧
            \item ARP提供IP地址到MAC地址的动态映射
        \end{itemize}
    \end{block}
    \vspace{0.3em}
    \begin{columns}[T]
        \column{0.48\textwidth}
        \textbf{ARP请求过程}
        \begin{enumerate}
            \item 主机A广播ARP请求："谁的IP是192.168.1.2？"
            \item 局域网内所有主机收到广播
            \item 目标主机B回复ARP响应："我是，我的MAC是AA:BB:CC:DD:EE:FF"
        \end{enumerate}
        \column{0.48\textwidth}
        \textbf{ARP缓存表（哈希表）}
        \begin{itemize}
            \item Key：IP地址
            \item Value：MAC地址 + TTL
            \item 典型超时：2-10分钟
            \item \textcolor{blue}{减少ARP广播次数}
        \end{itemize}
    \end{columns}
    \vspace{0.3em}
    \alert{易错点：ARP发生在IP路由\textbf{之后} --- 先确定下一跳IP，再通过ARP获取其MAC地址}
\end{frame}

\begin{frame}{数据封装成帧过程}
    \begin{center}
    \begin{tikzpicture}[>=stealth]
        \draw[fill=blue!10,rounded corners] (0,0) rectangle (10,0.6);
        \node at (5,0.3) {\small \textbf{应用层}：HTTP请求报文（GET /index.html）};
        \draw[fill=green!10,rounded corners] (0,0.8) rectangle (10,1.4);
        \node at (5,1.1) {\small \textbf{传输层}：TCP段（源端口+目的端口+序号+HTTP数据）};
        \draw[fill=yellow!10,rounded corners] (0,1.6) rectangle (10,2.2);
        \node at (5,1.9) {\small \textbf{网络层}：IP数据报（源IP+目的IP+TCP段）};
        \draw[fill=red!10,rounded corners] (0,2.4) rectangle (10,3.0);
        \node at (5,2.7) {\small \textbf{数据链路层}：帧（MAC+类型+IP数据报+CRC）};
        \draw[fill=gray!10,rounded corners] (0,3.2) rectangle (10,3.8);
        \node at (5,3.5) {\small \textbf{物理层}：比特流（010101...）};
        \draw[<->,thick] (10.5,0.3) -- (10.5,3.5);
        \node[right,font=\small] at (10.5,1.9) {每层添加头部};
    \end{tikzpicture}
    \end{center}
    \vspace{0.3em}
    \begin{itemize}
        \item 数据从高层向低层传递，每层添加自己的头信息
        \item 接收端逐层拆封，反向解析
    \end{itemize}
\end{frame}

\begin{frame}{计网：IP路由 --- ARP --- 封装成帧}
    \begin{block}{IP路由}
        \begin{itemize}
            \item 根据目的IP地址查询路由表，确定下一跳
            \item \textbf{路由表：Trie树（前缀树）} --- 最长前缀匹配算法
        \end{itemize}
    \end{block}
    \begin{block}{ARP协议}
        \begin{itemize}
            \item 已知IP地址，获取目标MAC地址（局域网内广播）
            \item ARP缓存表（哈希表）保存近期映射关系
        \end{itemize}
    \end{block}
    \begin{block}{封装成帧}
        \begin{itemize}
            \item 数据链路层添加帧头帧尾：MAC地址 + 类型 + CRC校验
            \item 物理层：比特流传输（曼彻斯特/4B/5B编码等）
        \end{itemize}
    \end{block}
\end{frame}

\begin{frame}{OS：Socket编程与内核交互}
    \begin{block}{Socket编程接口}
        \begin{itemize}
            \item 应用程序通过Socket API（\texttt{socket, bind, connect, send, recv}）进行网络通信
            \item Socket本质是操作系统提供的网络I/O接口
        \end{itemize}
    \end{block}
    \begin{block}{进程间通信与内核态切换}
        \begin{itemize}
            \item 用户态 $\to$ 内核态：系统调用（如 \texttt{send, recv}）触发软中断
            \item 内核协议栈处理TCP/IP协议，管理连接状态
        \end{itemize}
    \end{block}
    \begin{block}{数据缓冲区管理}
        \begin{itemize}
            \item 发送缓冲区、接收缓冲区（内核空间）
            \item \textbf{TCP缓冲区：队列结构} --- 字节流按序到达，等待应用程序读取
        \end{itemize}
    \end{block}
\end{frame}

\begin{frame}{计组：网卡DMA与中断处理}
    \begin{block}{网卡DMA传输}
        \begin{itemize}
            \item 网卡接收网络数据后，通过DMA直接写入内存缓冲区
            \item CPU无需逐字节搬运，提高效率
        \end{itemize}
    \end{block}
    \begin{block}{中断处理}
        \begin{itemize}
            \item 数据到达内存后，网卡发出中断信号
            \item CPU保存当前上下文，跳转到中断处理程序
            \item 中断处理完成后恢复用户进程
        \end{itemize}
    \end{block}
    \begin{block}{CPU执行协议栈指令}
        \begin{itemize}
            \item CPU执行内核协议栈代码（指令流水线）
            \item 频繁访问的代码和数据被Cache缓存，加速处理
        \end{itemize}
    \end{block}
\end{frame}

\begin{frame}{数据结构：四科交汇处的核心结构}
    \begin{block}{路由表 --- Trie树（前缀树）}
        \begin{itemize}
            \item 按IP前缀组织，支持最长前缀匹配
            \item 搜索复杂度 O(L)，L为IP地址长度
        \end{itemize}
    \end{block}
    \begin{block}{DNS缓存 --- 哈希表}
        \begin{itemize}
            \item Key：域名，Value：IP地址 + TTL
            \item 平均查找 O(1)，冲突处理（链地址法/开放定址法）
        \end{itemize}
    \end{block}
    \begin{block}{TCP缓冲区 --- 队列}
        \begin{itemize}
            \item 发送队列、接收队列，按序存储数据段
            \item 滑动窗口机制依赖于缓冲区的容量管理
        \end{itemize}
    \end{block}
\end{frame}

\begin{frame}{返回路径：服务器处理与HTTP响应}
    \begin{block}{服务器处理}
        \begin{itemize}
            \item 接收HTTP请求 $\to$ 解析请求行/首部 $\to$ 路由到处理程序
            \item 查询数据库（B+树索引）、读取文件，生成响应内容
        \end{itemize}
    \end{block}
    \begin{block}{HTTP响应返回}
        \begin{itemize}
            \item 响应报文：状态行 + 首部 + 空行 + 主体
            \item 再次经过TCP传输 $\to$ IP路由 $\to$ 链路传输
        \end{itemize}
    \end{block}
    \begin{block}{TCP四次挥手（连接关闭）}
        \begin{itemize}
            \item 数据传输完毕，发起FIN关闭连接（或保持Keep-Alive复用）
        \end{itemize}
    \end{block}
\end{frame}

\begin{frame}{浏览器渲染（核心流程）}
    \begin{enumerate}
        \item \textbf{HTML解析}：构建DOM树（文档对象模型）\textcolor{accent}{--- 树形结构}
        \item \textbf{CSS解析}：构建CSSOM树（CSS对象模型）
        \item \textbf{合并}：DOM + CSSOM $\to$ Render Tree（渲染树）
        \item \textbf{布局（Layout）}：计算每个节点的几何位置（盒模型）
        \item \textbf{绘制（Paint）}：将渲染树绘制到屏幕上（栅格化）
    \end{enumerate}
    \vspace{0.5em}
    \alert{浏览器渲染引擎（如Blink、WebKit）涉及栈、树等多种数据结构的综合运用}
\end{frame}

\begin{frame}{分层时间线图（完整流程）}
    \centering
    \begin{tikzpicture}[>=stealth, scale=0.85]
        % Draw layers
        \draw[fill=blue!5] (0,0) rectangle (13,0.5);
        \node[left] at (0,0.25) {\small 应用层};
        \draw[fill=green!5] (0,0.7) rectangle (13,1.2);
        \node[left] at (0,0.95) {\small 传输层};
        \draw[fill=yellow!5] (0,1.4) rectangle (13,1.9);
        \node[left] at (0,1.65) {\small 网络层};
        \draw[fill=red!5] (0,2.1) rectangle (13,2.6);
        \node[left] at (0,2.35) {\small 链路层};
        \draw[fill=gray!5] (0,2.8) rectangle (13,3.3);
        \node[left] at (0,3.05) {\small 物理层};
        % Timeline arrows
        \draw[->,thick,408blue] (1,0.25) -- node[above,font=\tiny] {HTTP} (2.5,0.25);
        \draw[->,thick,408blue] (2.5,0.95) -- node[above,font=\tiny] {TCP} (4,0.95);
        \draw[->,thick,408blue] (4,1.65) -- node[above,font=\tiny] {IP} (5.5,1.65);
        \draw[->,thick,408blue] (5.5,2.35) -- node[above,font=\tiny] {帧} (7,2.35);
        \draw[->,thick,408blue] (7,3.05) -- node[above,font=\tiny] {比特流} (8.5,3.05);
        % Reverse process (on server)
        \draw[->,thick,408orange] (9,3.05) -- node[below,font=\tiny] {接收} (10,2.35);
        \draw[->,thick,408orange] (10,2.35) -- (10.5,1.65);
        \draw[->,thick,408orange] (10.5,1.65) -- (11,0.95);
        \draw[->,thick,408orange] (11,0.95) -- (11.5,0.25);
        % Labels
        \node[below,font=\tiny] at (1.75,-0.3) {DNS/HTTP};
        \node[below,font=\tiny] at (3.25,-0.3) {三次握手};
        \node[below,font=\tiny] at (5.75,-0.3) {路由+ARP};
        \node[below,font=\tiny] at (10,-0.3) {服务器处理};
    \end{tikzpicture}
    \vspace{0.3em}
    \begin{block}{跨层理解}
        数据从高层逐层向下封装，经物理介质传输，到达目的地后逐层向上拆封。每层只关心自己的协议头，对上下层透明。
    \end{block}
\end{frame}

\begin{frame}{本节总结}
    \begin{enumerate}
        \item 输入URL到页面显示是一道经典的综合题，串联四门核心科目
        \item 计网负责传输：DNS $\to$ HTTP $\to$ TCP $\to$ IP $\to$ ARP $\to$ 帧
        \item OS提供Socket接口、进程通信与缓冲区管理
        \item 计组完成数据搬运：DMA传输 + 中断处理 + Cache加速
        \item 数据结构是底层骨架：哈希表(DNS缓存)、Trie树(路由表)、队列(TCP缓冲区)、树(DOM树)
    \end{enumerate}
\end{frame}

\begin{frame}{下节预告}
    \begin{center}
    \vspace{1em}
    {\Large \textbf{5-2 程序的一生}}\\
    \vspace{0.5em}
    \textcolor{408blue}{从源代码到进程运行的全过程}
    \vspace{1em}
    \begin{block}{下节内容预告}
        \begin{itemize}
            \item 编译器如何将源代码翻译为汇编指令？
            \item 链接器如何合并目标文件？
            \item 加载器如何将程序载入内存？
            \item 操作系统如何创建和管理进程？
        \end{itemize}
    \end{block}
    \vspace{1em}
    \textcolor{408orange}{敬请期待下一讲！}
    \end{center}
\end{frame}
\end{document}
